% The five hidden pictures: final double page (verso + facing recto).
% The pictures come from the 2026 textbook (its code/easter_egg_images.py), copied into figures/; they carry only
% the monograms. Never write here what the names mean to the family.
\newcommand{\oeuf}[4]{% #1 name, #2 file, #3 sentence, #4 expansion of the acronym
  \noindent{\sffamily\bfseries\LARGE\color{accent}\textls[180]{#1}}\hfill
    {\footnotesize\itshape\color{encre!70}#4}\par\vspace{4pt}
  \noindent\begin{tcolorbox}[enhanced,colback=white,colframe=encre!12,boxrule=0.5pt,arc=2pt,
      left=3pt,right=3pt,top=3pt,bottom=3pt,halign=center,drop fuzzy shadow=encre!25]
    \includegraphics[width=\linewidth,height=3.95cm,keepaspectratio]{#2}\end{tcolorbox}\par\vspace{2pt}
  \noindent{\small #3}\par}

\cleartoverso
\xglobal\colorlet{accent}{okorange}
\phantomsection
\addcontentsline{toc}{chapter}{The five hidden pictures}
\markboth{The five hidden pictures}{The five hidden pictures}
\thispagestyle{empty}
\vspace*{-0.6cm}
{\fontsize{25}{30}\selectfont The five hidden pictures\par}
\vspace{6pt}{\color{okorange}\rule{2.6cm}{1.4pt}}\par\vspace{8pt}
\noindent{\small\itshape The four letters of ALIX and the three of KAL hide in these theorems\ldots}\par
\vspace{4pt}
\noindent{\footnotesize\color{encre!80}A treasure hunt to finish with. Five theorems carry short names; each one
restates, under that name, a theorem already checked elsewhere by the judge, and all five are kept in the file
\texttt{Acronymes.lean}. The original theorems, replayed on 26 October, are level~A (for ALIX, subject to the upstream
openai/math theorem); the five short names are exact aliases of them, not replayed themselves. The pictures are level~X calculations, which illustrate a particular case without proving anything.}\par
\vspace{14pt}
\oeuf{ALIX}{egg_ALIX.pdf}{At any density, the triangular lattice has the smallest possible energy for every repulsive
force of the right family, if the upstream openai/math theorem is correct.}{A Lattice Is eXtremal}
\vspace{16pt}
\oeuf{KAL}{egg_KAL.pdf}{The chance that a mutation ends up taking over the whole population is a limit: the limit of
the probability of having been ``absorbed'' by the state in which everyone carries it, as the number of generations
grows.}{Kolmogorov Absorption Limit}
\clearpage

\thispagestyle{empty}
\vspace*{-0.6cm}
\oeuf{ILENE}{egg_ILENE.pdf}{On a network where everyone is replaced at the same rate, a mutant has the same chance of
winning wherever it appears: this is the isothermal law.}{Isothermal Law: Every Node Equivalent}
\vspace{9pt}
\oeuf{ALEXIS}{egg_ALEXIS.pdf}{A positive mixture of decaying exponentials keeps, for all its derivatives, the
alternating signs that make an admissible force: Laplace mixtures keep their signs.}%
{Averaged Laplace EXponentials Inherit Signs}
\vspace{9pt}
\oeuf{AURELIA}{egg_AURELIA.pdf}{Under selection, when each type keeps its fitness, the average fitness of the
population never decreases.}{Average Under Replication: Evolution Lifts It, Always}
\clearpage
